\chapter{Вычислительные эксперименты и анализ полученных результатов}

В данной главе описываются результаты вычисленных экспериментов, которые были
проведены с помощью разработанного программного обеспечения для моделирования
динамики опухолевых процессов: валидация математической модели на связанном
наборе клинических данных, параметрический анализ системы, а также
моделирование различных терапевтических стратегий и оценка их эффективности.

\section{Валидация математической модели на клинических данных}

Для оценки корректности используемой в работе математической модели была
проведена её валидация на основе клинических данных, полученных в эксперименте
\textit{Siu et al.} \cite{Siu1986}, в котором изучалось развитие B-клеточной
лимфомы ($BCL_1$) в селезенке испытуемых мышей. Эта выборка отличается наличием
выраженных колебательных процессов и эффекта насыщения системы, когда
увеличение количества опухолевых клеток не приводит к пропорциональному росту
иммунного ответа. Это дает возможность исследовать способность
модели воспроизводить нелинейные взаимодействия популяций, включая
фазы временной ремиссии и последующего рецидива заболевания.

Для обеспечения максимальной сходимости расчетной кривой динамики концентрации
опухолевых клеток во времени с данными клинических наблюдений была проведена
процедура глобальной оптимизации с использованием алгоритма дифференциальной
эволюции \cite{diff_evo}. Найденные алгоритмом оптимальные значения ключевых
коэффициентов отражены в таблице \ref{tab:optimized_params}.

\begin{table}[ht!]
  \centering
  \caption{Значения параметров модели, идентифицированные с помощью алгоритма
  дифференциальной эволюции.}
  \label{tab:optimized_params}
  \begin{tabular}{|c|c|p{6cm}|}
    \hline
    Параметр & Значение & Биологическая интерпретация \\
    \hline
    $c$ & $6.5 \times 10^{-7}$ & Скорость уничтожения опухолевых клеток \\
    \hline
    $d$ & $8.1 \times 10^{-9}$ & Скорость активации иммунных клеток \\
    \hline
  \end{tabular}
\end{table}

Результаты совмещения модельных данных численности опухолевых клеток и оцифрованных
точек клинических наблюдений изображены на рисунке \ref{fig:validation_plot}.
Статистический анализ точности аппроксимации показал, что коэффициент
детерминации $R^2$ равен $0.469$. Полученное значение указывает на способность
модели описывать порядка 47\% вариативности экспериментальных данных; вместе с
тем визуальный анализ траекторий подтверждает уверенное качественное совпадение
фаз роста и регрессии опухолевой популяции.

\begin{figure}[ht!]
  \centering
  \includesvg[width=0.8\textwidth]{fig/validation_plot.svg}
  \caption{Результаты валидации модели: сопоставление расчетной кривой и
  экспериментальных данных \textit{Siu et al.} \cite{Siu1986}.}
  \label{fig:validation_plot}
\end{figure}

Выявленные расхождения в частоте колебаний обусловлены тремя основными
допущениями, заложенными в структуру модели. Система не учитывает
пространственную распределенность и миграцию клеток. В реальном организме отток
лимфоцитов в лимфоузлы и процессы метастазирования служат естественным
демпфером, сглаживающим колебания численности популяций. Также в уравнениях
отсутствуют временные задержки. Модель игнорирует латентный период, необходимый
иммунной системе для распознавания антигена и деления клеток. Мгновенная
реакция системы на изменение концентраций искусственно учащает периоды
осцилляций в расчетах.

Тем не менее, использование алгоритмов глобальной оптимизации позволило успешно
калибровать систему для воспроизведения ключевых динамических режимов.
Обнаруженные расхождения не снижают прогностическую ценность разработанного программного
обеспечения при решении поставленных задач.

\section{Параметрический анализ и исследование чувствительности системы}

Для выявления ключевых параметров, оказывающих определяющее воздействие на
поведение модели, выполнен параметрический анализ системы. В первую очередь
оценивалось влияние параметра скорости пролиферации опухолевых клеток $a$ на
значение амплитуды колебаний концентрации опухолевых клеток. Это позволило
изучить роль фактора агрессивности опухоли в формировании общей
динамики системы (рисунок \ref{fig:a_amp}).

\begin{figure}[ht!]
  \centering
  \includesvg[width=0.95\textwidth]{fig/a_amp.svg}
  \caption{График зависимости амплитуды колебаний концентрации опухолевых клеток от
  параметра скорости пролиферации $a$.}
  \label{fig:a_amp}
\end{figure}

Полученная зависимость имеет отчетливую U-образную форму с локальным минимумом
в окрестности точки $a \approx 0.22$. Данная точка показывает, что в этой зоне
скорость пролиферации опухолевых клеток может быть компенсирована реакцией со
стороны иммунной системы организма. Вследствие этого амплитуда колебаний
концентрации опухолевых клеток оказывается минимальной.

Замечание о том, что даже пятикратное увеличение значения параметра $a$
приводит лишь к незначительному изменению амплитуды колебаний (до $25\%$),
указывает на неспособность варьирования параметра скорости пролиферации
опухолевых клеток оказать значительное влияние на предельную тяжесть протекания
заболевания. Этот параметр скорее влияет на частоту активизации рецидивов
опухоли, а максимальная концентрация опухолевых клеток в системе определяется
именно эффективностью их уничтожения, что определяет необходимость выявления
более критичных параметров, влияющих на борьбу с опухолевой популяцией.

Проведенный анализ чувствительности системы к изменению параметра $d$ (его
влияние на амплитуду колебаний концентрации опухолевых клеток) отражен на
рисунке \ref{fig:d_amp}.

\begin{figure}[ht!]
  \centering
  \includesvg[width=0.95\textwidth]{fig/d_amp.svg}
  \caption{Влияние параметра $d$ на максимальную концентрацию опухолевых клеток
  в системе.}
  \label{fig:d_amp}
\end{figure}

Область $d < 1 \cdot 10^{-8}$ указывает на диапазон значений, при которых система
достигает максимума концентрации опухолевых клеток ($10^8$ кл/мл) и его
достижение определяется изменением коэффициента скорости активации иммунных клеток.
Область значений $d \in [1 \cdot 10^{-8}; 3 \cdot 10^{-8}]$ характеризуется
снижением амплитуды колебаний -- система стабилизируется, переходя от
циклических колебаний к умеренным, амплитуда которых превышает начальный
уровень концентрации опухолевых клеток в системе не более чем на 1 порядок.
Эта область значений параметра $d$ указывает на способность системы к
поддержанию опухолевой нагрузки на организм на достаточно низком уровне при
отсутствии выраженных всплесков пролиферации опухолевых клеток. Такой вывод
указывает на высокую роль стимуляции иммунного ответа на общий исход течения заболевания.

Для количественной оценки накопленного патологического воздействия опухоли на
организм использовалась метрика $AUC$ (Area Under the Curve)
\cite{auc_metric_pharm}, рассчитываемая как определенный интеграл численности
популяции клеток по времени:
\begin{equation}
  AUC = \int_{0}^{t_{end}} T(t) dt.
\end{equation}

Эта величина описывает суммарную опухолевую нагрузку на организм за весь период
моделирования. В качестве ориентира летального порога послужило значение
$AUC_{limit} = 2.58 \cdot 10^{10}$ кл$\cdot$сут/мл, рассчитанное путем интегрирования
модельной кривой $T(t)$ на интервале от 0 до 90 суток в условиях отсутствия
терапии. Такой временной интервал соответствует клиническим данным из работы
\textit{Siu et al.} \cite{Siu1986}, согласно которым в контрольной группе мышей без
лечения к 90-м суткам фиксируется 100\% летального исхода.

Для сопоставления результатов в различных расчетных режимах использовалась
нормированная величина $AUC_{norm}$, определяющая кратность отклонения от
летального порога:
\begin{equation}
  AUC_{norm} = \frac{AUC_{param}}{AUC_{limit}}.
\end{equation}
Использование метрики $AUC_{norm}$ позволяет провести клиническую интерпретацию
полученных в результате моделирования результатов. Область значений $AUC_{norm}
\ge 1$ определяет критический уровень патологической нагрузки на организм,
достижение которого объясняется неспособностью организма перенести опухолевую
нагрузку, что в результате приводит к летальному исходу к концу моделируемого периода.

Анализ чувствительности параметров $a$ и $d$, проведенный совместно, позволил
выявить области их значений, при которых достигается взаимно противоположный
исход развития опухолевого процесса. Результат такого анализа отражен на
рисунке \ref{fig:heatmap}.

\begin{figure}[ht!]
  \centering
  \includesvg[width=0.5\textwidth]{fig/heatmap.svg}
  \caption{Тепловая карта зависимости параметра $d$ от значений параметра
  $a$ (цветовая шкала отражает значение интегральной опухолевой нагрузки).}
  \label{fig:heatmap}
\end{figure}

Полученная тепловая карта показывает разделение области критического исхода
течения опухолевого процесса и стабилизационной области, которые можно
заметить как горизонтальную границу на графике. Отмечается наличие четкой границы
на уровне значения параметра $d \approx 2 \cdot 10^{-9}$ мл/(клетка$\cdot$день).
При значениях свыше этой границы (верхняя часть графика) наблюдается
относительно небольшой уровень опухолевой нагрузки $AUC_{norm} \to 0$ на всём
диапазоне значений параметра $a$.

При этом нижняя часть графика ($d < 2 \cdot 10^{-9}$ мл/(клетка$\cdot$день))
показывает совершенно
иную опухолевую нагрузку. При сохранении темпов роста опухолевых клеток на
уровне $a < 3 \cdot 10^{-2}$ день$^{-1}$ на всем диапазоне значений параметра
$d$ наблюдается сравнительно небольшой уровень опухолевой нагрузки, что соответствует
области значений системы, при котором низкое значение темпов пролиферации
опухолевых клеток (параметр $a$) не влияет на летальность исхода прогрессирования опухоли.
В таком диапазоне значений параметра $a$ изменение параметра $d$
(скорость активации иммунных клеток) не имеет существенного значения с точки
зрения летальности исхода воздействия опухолевой популяции на организм.

Напротив, высокая скорость роста опухолевых клеток ($a \ge 3 \cdot 10^{-2}$ день$^{-1}$)
приводит систему к состоянию летального исхода, на что указывает красная
область на графике (рисунок \ref{fig:heatmap}), при котором нормированное
значение опухолевой нагрузки достигает значений $AUC_{norm} \ge 5$. В таком
случае, исправить исход протекания опухолевых процессов в организме способно
изменение параметра $d$ до значений свыше $d \ge 2 \cdot 10^{-9}$
мл/(клетка$\cdot$день) -- тогда
результирующий уровень нормированной опухолевой нагрузки сохраняется на
допустимом уровне, при котором протекание опухолевых процессов в организме не
приводит к летальному исходу. Этот факт свидетельствует о ключевом значении
параметра $d$ на исход борьбы с опухолью: его поддержание на уровне $d \ge 2
\cdot 10^{-9}$ мл/(клетка$\cdot$день) сохраняет допустимый уровень интегральной
опухолевой нагрузки на организм при варьировании параметра скорости
пролиферации опухолевых клеток в диапазоне $a \in [10^{-3}, 10^1]$.

С целью выявления взаимосвязи параметров скорости уничтожения опухолевых клеток
(параметр $c$) и скорости активации иммунных клеток (параметр $d$) был проведен
параметрический анализ с использованием метрики опухолевой нагрузки ($AUC$) на
организм на момент конца проведения процесса моделирования ($t_{end}$).
Результаты моделирования в виде графика тепловой карты представлены на рисунке
\ref{fig:heatmap_cd}.

\begin{figure}[ht!]
  \centering
  \includesvg[width=0.5\textwidth]{fig/heatmap_cd.svg}
  \caption{Тепловая карта зависимости интегральной нагрузки (AUC) от параметров
    эффективности цитотоксического уничтожения опухолевых клеток $c$ и
  мобилизации $d$. Видна доминирующая роль скорости рекрутирования.}
  \label{fig:heatmap_cd}
\end{figure}

Полученная тепловая карта позволяет уточнить условия <<ускользания>> опухоли
от иммунного надзора. Установлено наличие выраженной вертикальной границы
раздела фаз, проходящей по критическому значению параметра $d \approx 5 \cdot
10^{-10}$. В области ниже данного порога наблюдается режим неограниченного
прогрессирования, при котором даже многократное повышение коэффициента
цитотоксичности $c$ не приводит к снижению патологической нагрузки.

Данный результат имеет фундаментальное значение для обоснования приоритетных
направлений терапии. Он доказывает, что ключевым фактором выживания системы
является не индивидуальная агрессивность эффекторных клеток, а скорость и
массовость их мобилизации. Таким образом, стратегии, направленные на активацию
притока лимфоцитов (увеличение $d$), являются стратегически более оправданными,
чем попытки повышения их селективной токсичности (увеличение $c$) при
недостаточном темпе мобилизации.

Выявленные зависимости определяют дальнейшее определение стратегий терапии,
необходимой для подавление очага распространения опухолевых клеток. Результаты
параметрического анализа показывают, что для стабилизации системы и её перевода
в состояние устойчивой ремиссии критически важно повышение скорости активации
(параметр $d$) новых иммунных клеток в очаг опухоли.

\section{Моделирование терапевтических стратегий и оценка их эффективности}

В данном разделе описывается проведение вычислительных экспериментов с целью
выявления влияния различных сценариев терапии, при наличии и увеличении в
организме концентрации опухолевых клеток, на исход течения динамики опухолевых процессов.
Основной задачей является сравнительный анализ монотерапевтических и
комбинированных стратегий. Оценка эффективности каждой стратегии терапии основана на
анализе результатов проведенного моделирования динамики популяций иммунных и
опухолевых клеток в моделируемой системе.

\subsection{Анализ сценария терапии введения цитокинов}

Для изучения гипотезы о биологических ограничениях использования цитокинов в
качестве самостоятельного метода лечения был смоделирован сценарий монотерапии
интерлейкином-2 (IL-2). В ходе моделирования были установлены следующие параметры системы:
интенсивность введения цитокинов $s_C = 10^4$ нг/(мл·день), что соответствует
предельно допустимой терапевтической дозе \cite{cytokine_dose}, и коэффициент
$\lambda = 20.0$ день$^{-1}$, отражающий период полураспада свободных молекул
интерлейкина в организме. Коэффициент цитотоксического уничтожения опухолевых клеток был
зафиксирован на уровне $c = 6.0 \cdot 10^{-7}$, начальная численность популяции
$T_0 = 5 \cdot 10^5$ кл/мл.

Результаты численного моделирования динамики клеточных популяций в течение 120
дней представлены на рисунке \ref{fig:monotherapy}. На графике наблюдается
временная задержка при реакции системы на уровне 14 суток с момента начала терапии на 10-й
день моделирования.

\begin{figure}[ht!]
  \centering
  \includesvg[width=0.7\textwidth]{fig/cytokines_therapy.svg}
  \caption{Динамика численности опухолевых ($T$) и иммунных ($I$) клеток при
  монотерапии цитокинами ($s_C = 10^4, \lambda = 20$)}
  \label{fig:monotherapy}
\end{figure}

На 24-й день наблюдается незначительное повышение
концентрации иммунных клеток организма. Однако динамика опухолевой популяции не
показывает положительного исхода проведения терапии. На протяжении всего
временного диапазона проведения монотерапии цитокинами (с 10 по 50 день моделирования)
отмечается пиковая амплитуда концентрации опухолевых клеток на уровне $10^8$ кл/мл.
Такой эффект объясняется инерционностью системы: цитокинам
требуется некоторое время для усиления пролиферации эффекторных клеток до
уровня, способного замедлить увеличение популяции опухолевых клеток.
Но эффект снижения опухолевой популяции не достигается, о чем свидетельствует
анализ приведенного графика. Несмотря на значение опухолевой нагрузки
$AUC = 1.37 \cdot 10^9$ кл$\cdot$сут/мл ниже выявленного летального порога
($AUC_{limit} = 2.58 \cdot 10^{10}$ кл$\cdot$сут/мл) на конечный день
моделирования (120 день) введение высокой дозы
цитокинов ($s_C = 10^4$ нг/(мл·день)) объясняет нецелесообразность проведения
такой терапии в реальных клинических условиях, поскольку не приводит к
существенному улучшению динамики опухолевых процессов в организме.

\subsection{Изучение динамики системы при импульсной монотерапии адоптивным
переносом клеток}

При изучении поведения системы в условиях импульсной монотерапии адоптивным
переносом клеток рассматривался сценарий введения препарата, при котором
интенсивность притока эффекторных единиц $s_A$ составляла $40$
кл/(мл$\cdot$день) \cite{adoptive_therapy}. Такой сценарий предполагает
внешнее воздействие на систему в виде трех последовательных инъекций,
выполняемых на 20-е, 40-е и 60-е сутки процесса моделирования.

На начальном этапе, охватывающем период с 20 по 70 дни,
выбранная дозировка создает выраженный накопительный эффект (рисунок
\ref{fig:sa_monotherapy_analysis}). Это приводит к
стремительному падению концентрации опухолевых клеток с исходного высокого
уровня до экстремально низкого значения порядка $10^{-5}$ кл/мл.
Однако анализ кривой популяции иммунных клеток, выявляет эффект так называемого иммунного
провала. После прекращения инъекций на 60-й день моделирования уровень лимфоцитов
постепенно снижается, поскольку при достижении опухолью малых объемов
прекращается эффективная стимуляция притока новых эффекторных клеток, а внешняя
поддержка со стороны препарата полностью отсутствует.

В результате того, что в ходе терапии популяция опухолевых клеток не была
уничтожена до абсолютного биологического нуля, а иммунный надзор оказался
временно ослаблен, выжившая популяция опухолевых клеток показывает экспоненциальный
рост в интервале с 120 по 260 дни. К концу моделируемого периода опухоль достигает своего
пика в $10^8$ кл/мл и подавляет оставшиеся иммунные клетки организма.
Моделирование показывает, что визуальное исчезновение опухоли после 70 дня
является ложным, так как математическое равновесие системы остается смещенным в
сторону опухолевых клеток. Таким образом, монотерапия в объеме 3 инъекций по
40 кл/(мл·день) обеспечивает лишь временную ремиссию, что обосновывает
необходимость поиска
более интенсивных терапевтических стратегий.

\begin{figure}[htbp]
  \centering
  \includesvg[width=0.7\textwidth]{fig/adoptive_therapy_low.svg}
  \caption{Динамика концентрации клеток при проведении адоптивной терапии на
  20, 40 и 60 дни симуляции (доза $s_A =40$ кл/(мл·день))}
  \label{fig:sa_monotherapy_analysis}
\end{figure}

Для проверки гипотезы о существовании критической дозировки воздействия было
проведено моделирование сценария, в котором интенсивность притока эффекторных
единиц $s_A$ была увеличена до 70 кл/(мл·день) при сохранении прежнего
временного регламента инъекций на 20, 40 и 60 дни симуляции. Как показано на
рисунке \ref{fig:sa_elimination}, такая модификация сценария монотерапии
адоптивным переносом клеток приводит к качественному изменению поведения
системы по сравнению с предыдущим экспериментом.

\begin{figure}[ht!]
  \centering
  \includesvg[width=0.7\textwidth]{fig/adoptive_therapy_medium.svg}
  \caption{Динамика популяций клеток при увеличении дозировки $s_A =
  70$ кл/(мл·день) на 20, 40 и 60 дни симуляции}
  \label{fig:sa_elimination}
\end{figure}

В интервале между 40 и 60 днями после введения второй инъекции клеток
наблюдается необратимое снижение опухолевой популяции. Кривая
популяции $T$ падает ниже критический порог выживаемости и стремительно
устремляется к значениям ниже критического порога $10^{-9}$ кл/мл. В контексте
представленной математической модели данное значение рассматривается как
снижение количества опухолевых клеток до практически незначимых уровней,
где концентрация опухолевых клеток считается как пренебрежимо малый уровень концентрации.

Отсутствие признаков повторного всплеска концентрации опухолевых клеток на
промежутке моделирования в 300 дней указывает на достижение состояния устойчивой ремиссии.
Моделирование с увеличенной на 75\% дозой адоптивных клеток ($70$ кл/(мл·день))
указывает на наличие критического порога дозировки, после которой динамика
системы кардинально меняется. Такой результат определяет минимально необходимый
уровень дозировки адоптивных клеток, что может быть в дальнейшем использовано
при проведении комбинированной стратегии иммунотерапии для достижения
положительной динамики развития опухоли к концу расчетного периода.

С целью изучения влияния временного фактора на результаты адоптивной терапии
был проведен эксперимент с введением дозы адоптивных клеток величиной $s_A =
70$ кл/(мл·день), но уже начиная с 50 дня моделирования. Как видно на графике
(рисунок \ref{fig:sa_timing_window}) терапия из предыдущего эксперимента при
дозировке адоптивных клеток $s_A = 70$ кл/(мл·день) оказывается неэффективной
при её проведении, начиная с 50 дня моделирования. Динамика концентрации
опухолевых клеток при таком сценарии адоптивной терапии характеризуется высоким
уровнем осцилляций, при котором снижение концентрации опухолевых клеток
чередуется возвращением к их пиковой концентрации на уровне порядка $10^8$ кл/мл.

\begin{figure}[ht!]
  \centering
  \includesvg[width=0.7\textwidth]{fig/adoptive_therapy_timing.svg}
  \caption{Динамика системы при позднем начале терапии ($s_A = 70$ кл/(мл·день),
  начало терапии на 50-й день).}
  \label{fig:sa_timing_window}
\end{figure}

Результат моделирования такого сценария адоптивной терапии ($s_A = 70$
кл/(мл·день)), проводимой на 20 день моделирования, в сравнении с проведением
адоптивной терапии аналогичного уровня дозировки, но уже начиная с 50 дня,
показывает необходимость соблюдения временного интервала эффективности проводимой терапии.
Смещение окна начала терапии характеризуется кардинально противоположным эффектом динамики
концентрации опухолевых клеток. Этот факт указывает на необходимость учета момента начала
проведения терапии вместе с подбором её оптимальной дозировки для достижения
положительного эффекта адоптивной терапии.

\subsection{Анализ сценария комбинированной терапии адоптивным переносом клеток
и таргетными ингибиторами}

Для преодоления ограничений, выявленных при позднем начале монотерапии, был
изучен сценарий комбинированного воздействия, который совмещает адоптивный
перенос клеток с механизмами таргетного ингибирования \cite{target_therapy}. В
рамках данного эксперимента терапия проводилась начиная с 50-го дня моделирования
-- момента достижения пиковой концентрации опухолевых клеток на уровне порядка
$10^8$ кл/мл. Параметры комбинированного режима включали интенсивность притока
$s_A = 70$ кл/(мл·день), усиление коэффициента цитотоксичности до уровня
$\eta_c = 1.7 \cdot 10^{-6}$ и подавление интенсивности гибели иммунных клеток до
$\eta_{\mu} = 0.06$.

Результаты моделирования, представленные на рисунке \ref{fig:combined_50_day},
демонстрируют радикальное изменение динамики системы по сравнению с
монотерапией адоптивным переносом в идентичных текущему сценарию терапии
начальных условиях моделирования. Несмотря на начало проведения терапии
на этапе высокой концентрации опухолевых клеток клеток, синергетический
эффект выбранных механизмов терапии позволил остановить распространение,
обеспечив элиминацию опухолевых клеток.

\begin{figure}[ht!]
  \centering
  \includesvg[width=0.7\textwidth]{fig/combined_therapy.svg}
  \caption{Динамика комбинированной терапии при позднем начале терапии (50-й день):
  элиминация опухоли за счет синергетического эффекта $s_A$, $\eta_c$ и $\eta_{\mu}$}
  \label{fig:combined_50_day}
\end{figure}

За счет синергетического эффекта от проведения терапии адоптивным переносом иммунных
клеток ($s_A$), терапии усиления цитотоксичности иммунных клеток ($\eta_c$), а также
снижения гибели иммунных клеток под действием соответствующей терапии ($\eta_{\mu}$)
происходит качественное изменение динамики опухолевого процесса. Моделирование
показывает, что даже при условии начала терапии на 50-й день, когда
концентрация опухолевых клеток достигает своего пика, комбинированное
воздействие способно обеспечить полную элиминацию опухоли. Сравнивая
комбинированный сценарий проведенной терапии со сценарием монотерапии
адоптивным переносом, следует отметить, что комбинированная терапия
обеспечивает устойчивое подавление роста концентрации опухолевых клеток к 85
дню моделирования, тогда как монотерапия адоптивным переносом с тем же уровнем
дозировки ($s_A = 70$ кл/(мл·день)) и соблюдением идентичного момента начала
проведения терапии отмечается своей неэффективностью с точки зрения динамики
концентрации опухолевых клеток.

Проведенные сценарии терапии систематизированы в таблице \ref{tab:therapy_results}.

\begin{table}[htbp]
  \centering
  \caption{Сводные результаты моделирования и оценки эффективности
  терапевтических стратегий}
  \label{tab:therapy_results}
  \begin{tabular}{|p{4cm}|p{3cm}|p{3cm}|p{4cm}|}
    \hline
    Тип терапии & Параметры & AUC (120 дн.) & Клинический исход \\ \hline
    Без терапии & $s_A=0, s_C=0$ & $2.58 \cdot 10^{10}$ & Летальный исход к
    90-м суткам \\ \hline
    Монотерапия цитокинами & $s_C=10^4, \lambda=20$ & $1.37 \cdot 10^9$ &
    Отсутствие ответа, риск рецидива \\ \hline
    Адоптивная (малая доза) & $s_A=40$ (20, 40, 60 дн.) & $\approx 5.2 \cdot
    10^8$ & Ложное выздоровление, рецидив на 240 день \\ \hline
    Адоптивная (порог) & $s_A=70$ (20, 40, 60 дн.) & $\ll 1.0 \cdot 10^8$ &
    Полная элиминация, устойчивое выздоровление \\ \hline
    Адоптивная (поздний старт) & $s_A=70$ (старт на 50 день) & $1.15 \cdot
    10^9$ & Пропуск <<терапевтического окна>>, рецидив \\ \hline
    Комбинированная & $s_A=70, \eta_c, \eta_{\mu} > 0$ & $\ll 1.0 \cdot 10^8$ &
    Синергия, успех на запущенных стадиях \\ \hline
  \end{tabular}
\end{table}

Проведенные сценарии терапии адоптивным переносом указали на наличие корреляции
между началом её проведения и исходом результатом проведения терапии.
Установлено пороговое значение дозировки $s_A=70$ кл/(мл·день), которое
отмечается успешным исходом протекания опухолевых процессов в моделируемой системе
при соблюдении начала проведения терапии, соответствующего фазе роста
концентрации опухолевых клеток. Также сочетание адоптивной терапии вместе с
таргетным ингибированием (усиление цитотоксичности $\eta_c$ и снижение
гибели иммунных клеток $\eta_{\mu}$) показало синергетический эффект:
совместная терапия, проводимая начиная с 50 дня моделирования, привела к
успешному результату терапии, в ходе которой концентрация опухолевых клеток в
моделируемой системе на конечный день моделирования была зафиксирована на
уровне менее $10^{-9}$, что указывает на избавление от опухолевой популяции в
моделируемой системе.

Таким образом, разработанное программное обеспечение позволяет проводить предварительный
сравнительный анализ протоколов лечения, минимизируя риски неверного выбора
тактики иммунотерапии и обеспечивая подбор персонализированных стратегий на
основе математического прогноза.